% TOP_PROBLEM: {"problemId": "problem.rokhlin-problem-on-multiple-mixing", "canonicalTitle": "Rokhlin problem on multiple mixing", "releaseRank": 396, "primaryDomain": "probability_ergodic_dynamics", "primaryDomainLabel": "Probability, ergodic theory & dynamics", "displayStatus": "open", "releaseStatus": "open"}
% !TeX program = lualatex
% ATTEMPT_STATUS: unresolved
% Model: GPT-6 Astra Ultra; continuation dated 27 September 2026.
\documentclass[11pt]{article}
\usepackage[margin=25mm]{geometry}
\usepackage{fontspec}
\setmainfont{Latin Modern Roman}
\usepackage{amsmath,amssymb,mathtools}
\usepackage{unicode-math}
\setmathfont{Latin Modern Math}
\usepackage{xurl,hyperref}
\hypersetup{colorlinks=true,urlcolor=blue}
\setcounter{secnumdepth}{0}
\setlength{\parindent}{0pt}
\setlength{\parskip}{5pt}
\setlength{\emergencystretch}{3em}
\begin{document}
\section*{\texorpdfstring{Rokhlin problem on multiple mixing}{Rokhlin problem on multiple mixing}}\label{396}
\subsection{Definitions and mathematical statement}
Let $T$ be an invertible measure-preserving transformation of a probability space $(X,\mu)$. Mixing means $\mu(A\cap T^{-n}B)\to\mu(A)\mu(B)$ as $|n|\to\infty$ for all measurable $A,B$. The question asks whether this implies, for every $k\ge3$ and measurable $A_1,\ldots,A_k$,
\[
 \mu\!\left(A_1\cap T^{-n_1}A_2\cap\cdots\cap
 T^{-(n_1+\cdots+n_{k-1})}A_k\right)
 \longrightarrow\prod_{j=1}^k\mu(A_j)
\]
as all positive gaps $n_i\to\infty$. Invertibility is the convention compatible with the catalogue's two-sided mixing definition; a one-sided version must be stated separately. The assertion is mixing of all orders, not simply decay of two-point correlations for selected observables, and does not impose a quantitative rate of convergence.
\par\smallskip\textit{Formulation and concept references:} \href{https://arxiv.org/abs/2411.07234}{[S1]}.

\subsection{Short English statement}
Does ordinary mixing of a probability-preserving transformation force statistical independence of every finite collection of observations separated by large time gaps?
\subsection{Research attempt}
\paragraph{The order of the quantifiers.}
The inspected primary survey [S1], including its November 2024
revision, treats the implication from mixing to multiple mixing as
unresolved. The target concerns a single invertible transformation
and all fixed measurable test sets. It does not assume a mixing rate,
a finite-state Markov presentation, a Gaussian law, or a uniform
estimate over time-dependent observables. We prove the implication
under two additional structures and then isolate why neither argument
extends from the stated two-set hypothesis alone.

Write $t_1=0$ and $t_j=n_1+\cdots+n_{j-1}$. The desired limit is
equivalently
\[
 \int\prod_{j=1}^k f_j\circ T^{t_j}\,d\mu
       \longrightarrow\prod_{j=1}^k\int f_j\,d\mu
 \quad\hbox{as }\min_{i<j}|t_i-t_j|\longrightarrow\infty
 \tag{396.1}
\]
for bounded measurable $f_j$ at ordered times. Indicators give the
printed assertion; simple-function approximation gives the bounded
version. For increasing times, divergence of all consecutive gaps
is exactly divergence of the minimum pairwise distance. Invertibility
allows a simultaneous shift of all times without changing the integral.

\paragraph{A uniform approximation lemma.}
Suppose $|f_j|,|g_j|\leq1$. Telescoping the product gives pointwise
\[
 \left|\prod_j f_j\circ T^{t_j}-\prod_jg_j\circ T^{t_j}\right|
       \leq\sum_j|f_j-g_j|\circ T^{t_j}.
\]
Measure preservation and integration imply
\[
 \left|\int\prod_j f_j\circ T^{t_j}
       -\int\prod_jg_j\circ T^{t_j}\right|
       \leq\sum_j\|f_j-g_j\|_1.
 \tag{396.2}
\]
The same bound holds for the difference of the products of their
means. The crucial feature is independence of all times. Therefore
multiple mixing established on an $L^1$-dense collection of uniformly
bounded test functions extends to all bounded tests: fix an
approximation accuracy first, send the gaps to infinity second, and
only then send the approximation error to zero.

The lemma does not permit the approximating functions themselves to
depend on the large times without a uniform bound on their complexity.
That would change the order of quantifiers in the mixing assumption.

\paragraph{Independent coordinates: a complete all-orders subcase.}
Let $X=\mathcal A^{\mathbb Z}$ with product probability measure
$\nu^{\mathbb Z}$, and let $T$ be the two-sided coordinate shift.
Take bounded functions $g_j$ depending only on coordinates
$[-M,M]$. Their translates depend on windows
$[t_j-M,t_j+M]$. If every consecutive gap is greater than $2M$,
these windows are disjoint, hence independence gives the exact equality
\[
 \int\prod_jg_j\circ T^{t_j}\,d\mu
        =\prod_j\int g_j\,d\mu.
 \tag{396.3}
\]
Functions depending on finitely many coordinates are dense in $L^1$:
conditional expectations onto the increasing finite-coordinate
sigma-algebras converge in $L^1$ to any bounded measurable function,
and preserve its bound. Applying (396.2) proves mixing of every finite
order for this system, for arbitrary measurable tests. The proof
also covers arbitrary coordinate distributions; finite alphabets are
not needed for the independence step.

If a system is a measure-preserving factor of this shift, pull back
each test function along the factor map. Its multiple correlations
are the same as those of the pullbacks, so the same conclusion holds
for every such factor. This uses an actual factor map commuting with
the dynamics, not a similarity of two-point correlation functions.

\paragraph{A quantitative block condition sufficient for the conclusion.}
Consider more generally a stationary two-sided process $(X_j)$ that
generates the system. Define
\[
 \mathcal F_{\leq a}=\sigma(X_j:j\leq a),\qquad
 \mathcal F_{\geq b}=\sigma(X_j:j\geq b),
\]
and the process coefficient
\[
 \alpha(r)=\sup\{|\mu(A\cap B)-\mu(A)\mu(B)|:
      A\in\mathcal F_{\leq0},\ B\in\mathcal F_{\geq r}\}.
 \tag{396.4}
\]
For $[0,1]$-valued functions $F,G$ measurable with respect to the
two sigma-algebras, layer-cake integration gives
\[
 |\mathbb E(FG)-\mathbb EF\,\mathbb EG|\leq\alpha(r).
 \tag{396.5}
\]
Indeed express $F=\int_0^1\mathbf1_{F>s}\,ds$ and similarly for
$G$, apply (396.4) to every pair of level events and integrate.

Suppose $\alpha(r)\to0$. For $[0,1]$-valued cylinder tests depending
on $[-M,M]$, the product of the first $k-1$ shifted tests is measurable
before $t_{k-1}+M$, while the last is measurable after $t_k-M$.
Equation (396.5) separates the last test with error at most
$\alpha(t_k-t_{k-1}-2M)$. Iterating gives
\[
 \left|\mathbb E\prod_{j=1}^kg_j\circ T^{t_j}
       -\prod_{j=1}^k\mathbb Eg_j\right|
 \leq\sum_{j=1}^{k-1}\alpha(t_{j+1}-t_j-2M).
 \tag{396.6}
\]
Approximation by (396.2) now proves (396.1). Complex bounded tests
follow by splitting into real and imaginary parts and their positive
and negative parts. Thus this is a second explicit sufficient route.

It would be a serious mistake to identify (396.4) with the ordinary
mixing hypothesis. The latter fixes its two events before sending
the time to infinity. Equation (396.4) takes a supremum over large
past and future sigma-algebras. It is an extra condition on the chosen
process presentation. If the observation at each time recorded the
entire point of an invertible system, both sigma-algebras would already
be the whole measurable sigma-algebra, so this coefficient need not
decrease at all. A useful finite-coordinate observation is structure,
not something furnished by invertibility.

\paragraph{Gaussian systems: covariance controls every finite block.}
There is another complete subcase that does not use arbitrary
nonlinear independence from two-point correlations. Let a system be
generated by a centered stationary real Gaussian process $(X_j)$,
and assume its covariance $r(n)=\mathbb E(X_0X_n)$ tends to zero
as $|n|\to\infty$. Take any finite collection of coordinate windows
around the separated times $t_j$. The joint vector is Gaussian.
Its covariance matrix has fixed diagonal blocks and off-diagonal
entries $r(t_j-t_i+a-b)$, where $a,b$ range over fixed finite offsets.
Every off-diagonal block therefore tends to zero.

For a finite Gaussian vector with covariance $C$, its characteristic
function at $\xi$ is $\exp(-\xi^{\mathsf T}C\xi/2)$. The block
convergence just proved makes the limiting characteristic function
the product of those of the individual windows. It follows that the
joint window laws converge weakly to the independent product law.
This remains true when a diagonal covariance block is singular;
positive definiteness is not required in the characteristic-function
argument.

Thus correlations factor asymptotically for bounded continuous
functions of these windows. To obtain bounded measurable cylinder
tests, approximate each one in $L^1$ of its fixed marginal Gaussian
law by bounded continuous functions. The error is controlled by the
same telescoping argument as (396.2), uniformly in the separation,
because every marginal window distribution is fixed. Finally use
finite-coordinate conditional expectations to pass to all tests in
the generated sigma-algebra. This proves all-orders mixing for the
Gaussian system.

If the Gaussian system is assumed mixing in the original sense,
then $r(n)\to0$ follows as well. Approximate $X_0$ and $X_n$ by
bounded truncations; Cauchy--Schwarz and stationarity uniformly control
the truncation error by the fixed $L^2$ tails. Apply mixing to the
bounded truncations, then let the truncation level grow. Hence ordinary
mixing does imply all-orders mixing in this Gaussian subcase.
The decisive extra fact is that the full joint law is determined by
the covariance matrix. General processes do not have this property.

\paragraph{The invalid induction, displayed explicitly.}
At order three one would like to use mixing on
\[
 \int f\,(g\circ T^n)(h\circ T^{n+m})\,d\mu
 =\int f\,\bigl(g(h\circ T^m)\bigr)\circ T^n\,d\mu.
 \tag{396.7}
\]
For each fixed $m$, ordinary mixing does separate $f$ from the
parenthesized function as $n\to\infty$. The target sends both
$m,n\to\infty$, with no relation between their rates. The second
test in (396.7) varies with $m$, and the mixing assumption supplies
no uniformity on this family. Taking an iterated limit proves less
than the simultaneous-gap limit.

The lack of uniformity is unavoidable for the whole unit ball of
observables. If $A$ has measure strictly between zero and one, choose
$B_n=T^n A$. Then
\[
 \mu(A\cap T^{-n}B_n)-\mu(A)\mu(B_n)
       =\mu(A)(1-\mu(A))>0
\]
for every $n$. Thus no general theorem can upgrade ordinary mixing
to uniform mixing over all measurable pairs. A successful induction
would need special control of the particular moving products in
(396.7), not this false universal strengthening.

\paragraph{Pairwise independence does not identify a multiple limit.}
On the four equally likely triples
\[
 (0,0,0),\ (0,1,1),\ (1,0,1),\ (1,1,0),
\]
each pair of coordinates consists of independent fair bits, but the
probability of the all-zero triple is $1/4$, not $1/8$. Thus even
exact factorization of every two-coordinate marginal is insufficient
to force the three-coordinate product law.

There is a dynamical version of the same obstruction to a joining
argument. Take independent fair-bit Bernoulli sequences $x,y$ and
set $z_j=x_j+y_j\pmod2$ for every $j$. The law of $(x,y,z)$ is
invariant under the simultaneous shift and each pair of its three
sequence coordinates has product Bernoulli law. Nevertheless its
three-coordinate law is not the product, because $z=x+y$ holds
identically. Hence invariance plus pairwise product marginals does
not characterize a product self-joining.

This is not a counterexample to Rokhlin's question. The Bernoulli
shift was proved above to mix at every order, so this joining is
not a limit of its separated-time graph distributions of the relevant
kind. It demonstrates precisely the extra condition that a putative
joining proof would have to extract from being such a limit. Merely
checking the marginals and diagonal invariance leaves a real gap.

\paragraph{Diagnostics and outcome.}
The finite check enumerates the four parity triples, confirms every
pair marginal is uniform, and records the nonproduct triple marginal.
It also checks factorization for disjoint finite-coordinate Bernoulli
blocks. Such enumeration cannot test an infinite-time assertion for
an arbitrary transformation; it only audits the examples used to
reject an invalid proof route.

We have proved all-orders mixing for independent-coordinate shifts
and their factors, for process presentations satisfying the stated
past-future coefficient condition, and for mixing Gaussian systems.
For an arbitrary mixing transformation, the first unresolved step
is uniform control over the moving products in (396.7), or an
equivalent restriction that excludes nonproduct multiple limits.
None of the subcase proofs produces that control from the printed
hypotheses alone. The general question remains unresolved by this
attempt.
\subsection{Sources}
\begin{itemize}
\item[S1]V. V. Ryzhikov, 'Multiple mixing, 75 years of Rokhlin's problem', arXiv:2411.07234 (2024).
\url{https://arxiv.org/abs/2411.07234}.
\textit{Formulation source.}
\end{itemize}
\end{document}
